
LLM-as-judge has emerged as a practical alternative to execution-based evaluation for code correctness assessment. With execution requiring test infrastructure and compute resources, practitioners increasingly deploy language models to judge whether generated code will pass tests without running them. However, practitioners face a critical decision with limited empirical guidance: which model scale to deploy?

Common intuitions suggest that larger models provide higher accuracy and that combining multiple scales through ensemble voting should improve reliability. This work tests both assumptions through systematic scale-controlled comparison.

Prior work establishes that LLM code judges exhibit biases \citep{moon2025biases}, that proprietary models ''frequently misjudge'' correctness \citep{crupi2025effectiveness}, and that test-time scaling via MCTS can improve single-model accuracy \citep{wang2025mctsjudge}. However, no prior work has conducted scale-controlled comparison with fixed evaluation settings to isolate the effect of model scale on error patterns.

This study addresses this gap by investigating scale-dependent error patterns in LLM code judges. Four predictions were evaluated under standardized conditions (fixed zero-shot prompt, temperature=0, HumanEval+ ground truth):

\begin{enumerate}
\item \textbf{P1 (Scale ordering with diminishing returns)}: Judge-execution agreement increases with scale, but the $7B\rightarrow 70B$ gain is larger than $70B\rightarrow proprietary$.
\end{enumerate}

\begin{enumerate}
\item \textbf{P2 (Scale-dependent error patterns)}: Different scales exhibit statistically different FP/FN ratios.
\end{enumerate}

\begin{enumerate}
\item \textbf{P3 (Ensemble benefit)}: Scale-diverse majority voting outperforms the best single judge by $\geq 3$\%.
\end{enumerate}

\begin{enumerate}
\item \textbf{P4 (Unanimous reliability)}: When all scales agree, accuracy is $\geq 10$\% higher than when they disagree.
\end{enumerate}

Results confirm P1 and P2 with statistical significance (p=0.021 and p=3.$27\times 10^{-8}$ respectively), but falsify P3 and P4. Ensemble voting degrades accuracy by 8.05\% compared to the best single judge. Unanimous agreement correlates with lower accuracy than disagreement.

